\documentclass[11pt]{article}

\usepackage[margin=1in]{geometry}
\usepackage{amsmath}
\usepackage{booktabs}
\usepackage{hyperref}

\title{Road-testing and validating SynTangle:\\
a concise verification narrative}
\author{SynTangle development note}
\date{\today}

\begin{document}
\maketitle

\section{Purpose}

SynTangle is intended to produce minimally tangled multispecies chromosome
synteny layouts while preserving the biological order encoded by the input
assemblies.  Validation therefore has to establish more than the fact that a
drawing looks cleaner.  It must show that the program:

\begin{enumerate}
  \item never uses an illegal subchromosomal transformation;
  \item finds or bounds the best legal whole-chromosome layout;
  \item is insensitive to arbitrary presentation scrambling;
  \item behaves sensibly on known simulated chromosome histories;
  \item scales in a measurable way as biological and combinatorial complexity
        increase; and
  \item produces outputs that remain interpretable by direct visual inspection.
\end{enumerate}

This document gives the minimum conceptual verification sequence needed to
make those claims.  It intentionally omits many implementation-level
experiments and debugging steps.

\section{What is being validated}

The biological input is naturally n-ary.  A homologous block may occur in
chromosomes from many species at once, so the conceptual representation is a
multispecies hypergraph.  SynTangle operates on an exact bipartite incidence
representation of that hypergraph: chromosome vertices on one side,
homology-group vertices on the other, and one incidence edge for each block
occurrence.

The legal display state for a species consists only of a permutation of whole
chromosomes and whole-chromosome orientation choices.  Internal chromosome
order is immutable.  The central optimization target is the crossing count
under those legal transformations.

The strongest result for a case is therefore not merely a low crossing count
but a statement of the form
\[
  C_{\mathrm{LB}} = C_{\mathrm{UB}} = C^*,
\]
where \(C^*\) is the constrained optimum.  When equality cannot be proved,
SynTangle must retain and report the nonzero optimality gap rather than
silently treating the best layout found as optimal.

\section{Verification architecture}

The minimum simulation-based validation separates biological evolution from
presentation noise:

\begin{verbatim}
known ancestor
    |
    v
logged fusion / fission / inversion history
    |
    v
extant simulated genomes
    |
    v
hidden native presentation
    |
    v
whole-chromosome permutation / reversal only
    |
    v
PUBLIC tangled benchmark input
    |
    v
SynTangle
    |
    v
optimized legal layout
    |
    v
only then reveal hidden simulator information
\end{verbatim}

The hidden evolutionary log and hidden tangle log are retained for validation
but are never supplied to the solver.  This separation is essential: otherwise
a successful benchmark could simply reflect leakage of the simulated answer.

\section{Step 1: invariant and tiny exact tests}

Before testing scale, the legal state space must be tested on deliberately
small fixtures whose behavior can be understood exactly.

The essential cases are:

\begin{itemize}
  \item perfect one-to-one chromosome correspondence;
  \item presentation scrambling with no biological change;
  \item a case requiring whole-chromosome reversal;
  \item a case in which an apparently cleaner solution would require an
        illegal subchromosomal edit;
  \item simple fusion/fission structures;
  \item balanced orientation constraints;
  \item contradictory orientation constraints, which must be reported rather
        than silently coerced; and
  \item small cyclic cases with a nonzero intrinsic crossing optimum.
\end{itemize}

These tests establish that the solver is searching the intended legal problem,
rather than merely minimizing crossings in an unconstrained drawing.

\section{Step 2: compare independent exact formulations}

For sufficiently small cases, the same optimum should be recoverable by more
than one exact method.  In SynTangle this includes direct/tiny exact search,
species-layer dynamic programming, and bounded branch-and-bound when allowed
to exhaust its search.

Agreement among independent exact formulations is useful because the methods
make different computational assumptions while solving the same legal
objective.  A disagreement is a correctness failure, not a performance issue.

\section{Step 3: presentation-invariance benchmark}

A key property of the problem is that arbitrary display order should not
change the constrained biological optimum.

For each simulated biological dataset, create three public inputs from the
same extant chromosomes:

\begin{center}
\begin{tabular}{lll}
\toprule
Presentation & Biological evidence & Display scrambling \\
\midrule
mild   & identical & low \\
strong & identical & high \\
random & identical & random order/orientation \\
\bottomrule
\end{tabular}
\end{center}

The biological fingerprint, hidden native crossing baseline, and solver seed
are held fixed within the triplet.  If all three runs are proved optimal, they
must satisfy
\[
 C^*_{\mathrm{mild}}
 =
 C^*_{\mathrm{strong}}
 =
 C^*_{\mathrm{random}}.
\]

This is a stronger test than asking whether each run removes many crossings:
it verifies that SynTangle is recovering a property of the biological problem
rather than a property of the starting drawing.

\section{Step 4: visual ground truth}

Numerical validation is not sufficient for a layout program.  Representative
cases should be inspected as a three-stage visual audit:

\begin{verbatim}
hidden simulator-native layout
        |
        v
deliberately tangled public layout
        |
        v
SynTangle optimized layout
\end{verbatim}

The visual report should also show the raw incidence structure, structural
projection diagnostics, crossing counts, and proof status.

The purpose of this inspection is not to demand that the optimized layout
reproduce the simulator-native drawing exactly.  The simulator-native
presentation is only one legal baseline and need not be the unique optimum.
Instead, visual inspection should confirm that:

\begin{itemize}
  \item chromosomes move or reverse only as whole units;
  \item internal genomic order is preserved;
  \item obvious presentation tangles are removed;
  \item residual crossings correspond to genuine constraints rather than
        unexplained display artifacts; and
  \item increasing benchmark complexity produces biologically plausible,
        inspectable outputs rather than pathological layouts.
\end{itemize}

\section{Step 5: controlled scaling}

Easy cases are useful for correctness but not for identifying the computational
limits of the method.  Complexity should therefore be increased along
separable and then coupled axes.

The principal axes are:

\begin{enumerate}
  \item number of species;
  \item number of chromosomes;
  \item number and combination of fusion/fission/inversion events; and
  \item anchor density, as a secondary evidence-size axis.
\end{enumerate}

The current stress design deliberately increases the first three together.
The aim is not to force every large case to prove an optimum.  The first
reproducible bounded case is valuable because it identifies the point at which
the current solver ceases to close the lower/upper-bound gap efficiently.

For every run, retain at least:

\begin{itemize}
  \item input crossing count;
  \item best legal crossing count found;
  \item lower bound and gap;
  \item proof status;
  \item solver selected;
  \item states or nodes evaluated;
  \item wall time;
  \item raw and structural graph summaries; and
  \item residual decision-space summaries.
\end{itemize}

\section{Step 6: explain computational difficulty using the residual problem}

Raw chromosome count alone is not expected to determine runtime.  After exact
reductions, SynTangle constructs a residual variable/factor description whose
variables are unresolved whole-chromosome order and GF(2) orientation choices,
and whose factors are crossing terms.

This representation permits more meaningful complexity diagnostics, including:

\begin{itemize}
  \item number and size of disconnected residual factors;
  \item objective-neutral unresolved variables;
  \item articulation variables;
  \item a heuristic upper bound on residual treewidth;
  \item implicit orientation-state count;
  \item orientation and order nodes actually evaluated; and
  \item memoized subproblem reuse.
\end{itemize}

The central scaling question is therefore whether runtime is better predicted
by the size/width of the residual coupled problem than by the size of the
original genome representation.

\section{Step 7: verify decomposition and parallelization}

Independent chromosome--homology incidence components are exact additive
factors of the current objective.  Serial and parallel solutions of those
components must therefore produce identical global lower bounds, upper bounds,
and optimized crossing counts.

Finer decomposition is permitted only when the residual objective is also
proved to factor.  The intended hierarchy is:

\begin{verbatim}
independent benchmark cases
    ->
independent incidence components
    ->
disconnected residual factor components
    ->
small-separator / tree-decomposition subproblems
    ->
parallel branch-and-bound frontier
\end{verbatim}

Parallelism must never change the mathematical result or weaken proof
bookkeeping.  A global optimum is proved only after every independent or
conditional subproblem is exhausted or bounded above the final incumbent.

\section{Step 8: ablation of algorithmic stages}

A method should not remain in the claimed core simply because it is
mathematically attractive.  Important stages should eventually be ablated on
the same fixed benchmark cases.

For each candidate stage, compare the solver with and without it while holding
simulation and presentation seeds fixed.  Record changes in:

\begin{itemize}
  \item correctness and proof status;
  \item lower-bound strength;
  \item states/nodes evaluated;
  \item wall time; and
  \item interpretability or audit value.
\end{itemize}

This is particularly important for structural projection, bridges,
articulation points, 2-cores, biconnected structure, cycle bases, and future
separator/treewidth machinery.  These should be described as active solver
machinery only after demonstrating a computational or proof benefit.

\section{Step 9: real-data road test}

Simulation establishes correctness because hidden truth is available, but it
does not establish usefulness on biological data.

A later road test should therefore use one or more real multispecies synteny
datasets with:

\begin{itemize}
  \item chromosome-scale assemblies;
  \item a transparent homology/synteny-block construction;
  \item known or interpretable fusion/fission relationships where possible;
  \item a sufficiently small subset for direct visual inspection; and
  \item a larger subset that exercises the scaling machinery.
\end{itemize}

The real-data test should ask whether SynTangle produces a cleaner legal
representation, whether residual tangles have a biological explanation, and
whether the result is stable to deliberate scrambling of the initial
chromosome display order.

No ancestral-event claim is required for this primary validation.

\section{Step 10: later ancestral-history validation}

Ancestral fusion/fission/inversion inference is a separate downstream problem.
When implemented, candidate histories must be constrained by a species tree
and evaluated against hidden simulator histories.

Because multiple histories may explain the same extant structures, validation
should measure ancestral adjacency recovery, event/breakpoint error, and the
size of the equivalence class of equally compatible histories rather than
requiring exact recovery of one hidden event sequence.

\section{Minimum evidence package for a methods claim}

A concise validation package for a paper or release should contain:

\begin{enumerate}
  \item invariant/tiny-fixture tests showing that illegal transformations are
        impossible;
  \item independent exact-method agreement on small cases;
  \item paired presentation-invariance benchmarks;
  \item native/tangled/optimized visual audits;
  \item a scaling ladder that reaches at least one genuinely difficult regime;
  \item explicit lower/upper bounds and proof status;
  \item serial-versus-parallel equivalence tests for exact decompositions;
  \item ablation evidence for any graph/decomposition stage claimed to improve
        the solver; and
  \item at least one real-data demonstration.
\end{enumerate}

This sequence is intentionally shorter than the complete development history.
It is the conceptual chain needed to support the claims that SynTangle solves
the correct constrained problem, removes presentation noise, reports intrinsic
residual tangledness honestly, and scales by exploiting mathematical structure
rather than by relaxing the biology.

\end{document}
